\documentclass[10pt,a4paper]{amsart}
\usepackage[margin=19mm]{geometry}
\usepackage{amsmath,amssymb,mathtools}
\usepackage[scr=rsfs,cal=euler]{mathalfa}
\usepackage{xfrac}
\usepackage[unicode,hidelinks,bookmarks=false]{hyperref}
\newcommand{\ie}{i.e.\ }
\newcommand{\cf}{cf.\ }
\newcommand{\resp}{respectively\ }
\newcommand{\Subsection}{\subsection}
\providecommand{\nobreakdash}{\nobreak\hbox{-}}
\setlength{\emergencystretch}{2em}
\multlinegap0pt
\usepackage{bm}
\usepackage{bbm}
\usepackage{dsfont}
\usepackage{xspace}
\usepackage{cancel}
\usepackage{mathtools}
\usepackage{amsthm}
\makeatletter
\@ifundefined{thm}{\newtheorem{thm}{Thm}}{}
\@ifundefined{lem}{\newtheorem{lem}[thm]{Lemma}}{}
\makeatother
\renewcommand{\tilde}{\widetilde}
\title[Moduli spaces of curves with polynomial point counts]{Moduli spaces of curves with polynomial point counts}
\author{Samir Canning, Hannah Larson, Sam Payne, Thomas Willwacher}
\date{Source: arXiv:2410.19913v2 (CC BY 4.0)}
\begin{document}
\maketitle

\noindent\emph{Source and licence.} Moduli spaces of curves with polynomial point counts. Version arXiv:2410.19913v2 (CC BY 4.0), distributed under the Creative Commons Attribution 4.0 International licence, as recorded in the arXiv licence metadata for this exact version and shown on the abstract page https://arxiv.org/abs/2410.19913v2. Full rights notice: licenses/arxiv-2410.19913-CC-BY-4.0.txt.\par\smallskip

\noindent\textit{Editorial scope.} Counted unit: the Lem whose source label is \texttt{lem:Deltap mono} in the article (source0.tex lines 2448--2456, source label \texttt{lem:Deltap mono}), delivered together with the article's own complete original proof of it (source0.tex lines 2457--2525, one \texttt{proof} environment of 69 lines). No mathematics is written, repaired or re-derived: statement and proof are the article's own text line for line. Required context retained uncounted: lem:phibd (lines 2354--2359). Every definition, display and label of those lines is retained unchanged. Omitted from this package and not redistributed: 1--2353, 2360--2447, 2526--3751. Nothing inside a retained excerpt is omitted or altered: every retained body is delivered exactly as the article's lines are written, so no omission receipt of any kind accompanies this package. Citations: the retained excerpts carry no citation command at all, so no bibliography entry of the article is transcribed and no reference is redistributed. Reference adaptation: none was needed and none was made; every reference inside the delivered unit resolves inside it, so no unresolved reference or citation is produced. Printed slips kept literal: no printed word, symbol, hypothesis or number of the counted statement or of its proof was altered. Layout and class: the article's own class is \texttt{amsart} with packages including tikz-cd, amssymb, amsmath, amscd, mathrsfs, url, cite, fullpage, hyperref, verbatim, comment, fancyvrb, fvextra, caption; the wrapper is the lane's pinned \texttt{amsart} editorial wrapper at 10pt on A4, which restates the theorem-like environments the retained text needs and adds the article's own macro definitions that the retained text needs (\texttt{\textbackslash tilde}), with extra wrapper packages (bm, bbm, dsfont, xspace, cancel, mathtools). The delivered numbering is the unit's own. Assets: no retained span contains a figure environment, a graphics-inclusion command, a PDF-page inclusion, a table or a TikZ picture; no image file is copied into this package, the package declares no asset dependency and no image file is redistributed with it. Licence: version arXiv:2410.19913v2 of this article is distributed under the Creative Commons Attribution 4.0 International licence, as recorded in the arXiv licence metadata for this exact version and shown on the abstract page https://arxiv.org/abs/2410.19913v2. A full rights notice is delivered at licenses/arxiv-2410.19913-CC-BY-4.0.txt. Characters: every retained excerpt is pure ASCII, so the delivered engine, XeLaTeX with its default Latin Modern text font, reports no missing character.
\begin{lem} \label{lem:phibd}
For any nonnegative real $x$, we have
 \begin{align*}
    \frac{x^{\Gamma+1}}{(\Gamma+1)!} \leq \phi_\Gamma(x)&\leq \frac{1}{\Gamma!} x^\Gamma e^x 
\end{align*}   
\end{lem}

\begin{lem}\label{lem:Deltap mono}
    Let $\Gamma\leq 10$ and $5\geq k_0\geq 2$.
Then the sequence $\Delta_{g,k_0,\Gamma}$ is monotonically decreasing in $g$ for $g\geq 150$. 

Furthermore, 
\[
\lim_{g\to \infty} \Delta_{g,k_0,\Gamma}=0.
\]
\end{lem}

\begin{proof}
We may show both statements at once by showing that the sequence
\[
    \tilde\Delta_{g,k_0}:=(g-2)\Delta_{g,k_0,\Gamma}
\]
is monotonically decreasing for $g\geq 150$. Here we suppress $\Gamma$ from the notation.
To this end, let $a_k=(3.1)^{k-1}\frac{(2\zeta(2))^k }{(2\pi)^{k-1}k!}\phi(2\pi k)$. For fixed $k$, the summands $a_k \frac{(g-2k)!}{(g-3)!}$ are monotonically decreasing in $g$.
Hence, if $g$ is odd, then clearly 
\[
    \tilde\Delta_{g+1,k_0}\leq \tilde\Delta_{g,k_0}.
\]
For $g$ even, there is one more summand in the expression for $\tilde\Delta_{g+1,k_0}$, and we can estimate 
\[
    \tilde\Delta_{g+1,k_0}-\tilde\Delta_{g,k_0}
    \leq  a_{k_0}\left(\frac{(g+1-2k_0)!}{(g-2)!} -\frac{(g-2k_0)!}{(g-3)!}\right) + \frac{a_{g/2}}{(g-3)!}.
\]
Hence, our sequence is monotonically decreasing if 
\begin{equation}\label{eq:Condition}
    a_{g/2}
    \leq 
    a_{k_0}
    (g-2k_0)!\left(1-\frac{g+1-2k_0}{g-2}\right).
\end{equation}
Inserting the definition of $a_k$, the inequality \eqref{eq:Condition} becomes
\[
 \left(3.1\frac{2\zeta(2)}{2\pi}\right)^{g/2-k_0}
\frac{\phi_\Gamma(\pi g)}{\phi_\Gamma(2\pi k_0)}
\leq (g-2k_0)! \frac{(g/2)!}{k_0!}\frac{2k_0-3}{g-2}.
\]
Applying Lemma \ref{lem:phibd}, we see that
inequality \eqref{eq:Condition} is implied by 
\[
    (\Gamma+1) \left(3.1\frac{2\zeta(2)}{2\pi}\right)^{g/2-k_0}
    \frac{(\pi g)^\Gamma e^{\pi g}}{(2\pi k_0)^{\Gamma+1}}
    \leq (g-2k_0)! \frac{(g/2)!}{k_0!}\frac{2k_0-3}{g-2}.
\]
Note that $\pi+\log(3.1\frac{2\zeta(2)}{2\pi})/2\leq 4$ so that
\[
    \left(3.1\frac{2\zeta(2)}{2\pi}\right)^{g/2}e^{\pi g} \leq e^{4g}.
\]
Hence, it suffices to show
\[
    \left(3.1\frac{2\zeta(2)}{2\pi}\right)^{-k_0} \frac{(\Gamma+1)}{\pi(2 k_0)^{\Gamma+1}}
    g^\Gamma e^{4 g}
    \leq (g-2k_0)! \frac{(g/2)!}{k_0!}\frac{2k_0-3}{g-2}.
\]

Furthermore, $g! \geq e(g/e)^g$ and $(g-2k_0)!>g! g^{-2k_0}$.
Therefore, \eqref{eq:Condition} is implied by 
\[
    \left(3.1\frac{2\zeta(2)}{2\pi}\right)^{-k_0}    
    \frac{(\Gamma+1)k_0!}{\pi(2k_0-3)(2k_0)^{\Gamma+1}}
    \leq e^{-5g+1} g^g
    \frac{(g/2)!}{g^{\Gamma+2k_0+1}}.
\]
The constant on the left-hand side is, for $k_0=2,\dots, 5$ and $\Gamma=0,\dots, 10$ bounded above by
\[
    \left(3.1\frac{2\zeta(2)}{2\pi}\right)^{-2}    
    \frac{(10+1)5!}{\pi(4-3)(4)^{0+1}}\lessapprox 40.
\]
Suppose that $g\geq e^5\approx 148.413$. Then $e^{-5g+1} g^g\geq 1$. For $k_0\leq 5$ and $\Gamma\leq 10$ 
inequality \eqref{eq:Condition} is hence implied by
\[
    40 
    \leq 
    \frac{(g/2)!}{g^{21}}.
\]
If $g\geq 100$ then $\frac{(g/2)!}{g^{21}}\gg 40$, and thus the inequality \eqref{eq:Condition} is satisfied. 
\end{proof}

\end{document}
